\chapter{Worked Examples and Exam Map}
\section{How the Topics Map to Exam-Style Thinking}
\begin{tabularx}{\textwidth}{@{}L{0.24\textwidth}Y@{}}
\toprule
Topic & What to recognize \\
\midrule
Units & Convert prefixes, distinguish \(f\) from \(\omega\), and check dimensions. \\
Complex impedance & Use rectangular form for series addition and polar form for magnitude/phase. \\
Resonance & Ask series or parallel, then ask whether the question concerns current, impedance, voltage, or bandwidth. \\
Filters & Identify passband, stopband, cutoff, slope, ripple, and bandwidth. \\
Decibels & Use \(10\log\) for power and \(20\log\) for voltage or current ratios under equal impedance. \\
Transmission lines & Use \(\Gamma\), SWR, return loss, velocity factor, and wavelength. \\
Active circuits & Separate gain, feedback, bandwidth, distortion, and stability. \\
\bottomrule
\end{tabularx}

\section{Example: Capacitive Reactance}
Find \(X_C\) for \(C=\SI{100}{\pico\farad}\) at \(f=\SI{14.2}{\mega\hertz}\).
\[
\omega=2\pi f=2\pi(14.2\times10^6).
\]
\[
X_C=\frac{1}{\omega C}
=\frac{1}{2\pi(14.2\times10^6)(100\times10^{-12})}
\approx \SI{112}{\ohm}.
\]
The impedance is \(Z_C=-\jj\SI{112}{\ohm}\).

\section{Example: Series RLC Resonance}
Given \(L=\SI{2.2}{\micro\henry}\) and \(C=\SI{150}{\pico\farad}\),
\[
f_0=\frac{1}{2\pi\sqrt{LC}}
\approx \SI{8.76}{\mega\hertz}.
\]
If \(R=\SI{3}{\ohm}\),
\[
Q_s=\frac{2\pi f_0 L}{R}
\approx 40.4,
\qquad
BW\approx \frac{8.76\times10^6}{40.4}
\approx \SI{217}{\kilo\hertz}.
\]

\section{Example: Decibel Ratio}
A filter attenuates a voltage by a factor of 0.25 with the same source and load
impedance. The voltage ratio in decibels is
\[
20\log_{10}(0.25)= -\SI{12.0}{\decibel}.
\]
The corresponding power ratio is \(0.25^2=0.0625\), which is also
\[
10\log_{10}(0.0625)= -\SI{12.0}{\decibel}.
\]

\section{Example: SWR}
If \(|\Gamma|=0.2\), then
\[
\mathrm{SWR}=\frac{1+0.2}{1-0.2}=1.5.
\]
The reflected power fraction is \(|\Gamma|^2=0.04\), or \(4\%\).

\section{Final Study Pattern}
Start every circuit problem with conservation. Ask what stores energy, what
dissipates energy, and whether the problem is best read as time-domain,
state-space, phasor, frequency-response, or transmission-line behavior. That
habit turns the manuscript's central idea into a practical exam strategy.
